\chapter{Nasal Congestion (Blocked Nose)}

\section*{Definition}
Obstruction of nasal airflow due to inflammation, swelling of nasal turbinates, mucus accumulation, or structural abnormalities, causing difficulty breathing through the nose.

\section*{What Happens in Your Body}
Nasal congestion results from widening of blood vessels and engorgement of venous sinusoids in the nasal turbinates triggered by inflammation (allergic, infectious), autonomic dysregulation, or medication effects. Mucosal edema narrows nasal passages while increased secretions further obstruct airflow. Chronic causes include structural abnormalities (deviated septum, nasal polyps) or chronic rhinosinusitis.

\section*{Causes}
\begin{itemize}
  \item Common cold (viral upper respiratory infection)
  \item Allergic rhinitis (hay fever)
  \item Sinusitis (acute or chronic)
  \item Deviated nasal septum
  \item Nasal polyps
  \item Vasomotor rhinitis
  \item Pregnancy rhinitis
  \item Rhinitis medicamentosa (decongestant spray overuse)
  \item Environmental irritants (smoke, pollution, dry air)
  \item Enlarged adenoids or tonsils (in children)
\end{itemize}

\section*{When to See a Doctor}
See your doctor if:
\begin{itemize}
  \item Symptoms are severe or getting worse over time
  \item Nasal congestion with high fever, facial swelling, visual changes, or unilateral congestion with bloody discharge
  \item You have other concerning symptoms such as fever, weight loss, or bleeding
\end{itemize}

\section*{Related Symptoms}
\begin{itemize}
  \item Runny nose
  \item Sneezing
  \item Sinus pressure
  \item Headache
  \item Reduced sense of smell
\end{itemize}
\textbf{Important:} If this symptom persists, your doctor will check for serious underlying causes.

\section*{Self-Care / Home Management}
\begin{itemize}
  \item Rest and avoid activities that make it worse.
  \item Maintain adequate hydration and a balanced diet.
  \item Over-the-counter remedies may provide symptomatic relief (consult a pharmacist or doctor).
  \item Monitor symptoms and seek professional advice if they persist or worsen.
  \item Avoid self-medicating with prescription medications without medical guidance.
\end{itemize}

\section*{How Your Doctor Will Evaluate You}
\begin{itemize}
  \item A thorough history and physical examination are the first steps in evaluation.
  \item Laboratory tests (blood work, urinalysis) may be ordered based on clinical suspicion.
  \item Imaging studies (X-ray, ultrasound, CT, MRI) may be indicated when structural pathology is suspected.
  \item Specialist referral is arranged when the diagnosis remains uncertain or requires specific expertise.
\end{itemize}

\section*{Treatment Options}
\begin{itemize}
  \item Treatment targets the underlying cause identified through diagnostic evaluation.
  \item Supportive care and symptom management are provided while awaiting definitive therapy.
  \item Medications, lifestyle modifications, and physical therapies may be prescribed as appropriate.
  \item Follow-up ensures treatment effectiveness and allows timely adjustments to the care plan.
\end{itemize}

\clearpage